\textbf{Время:} резерв

\textbf{Тезис:} современный метод соединяет общую квантовую основу и развитие способов пространственно-временного анализа.

\textbf{Визуал:} исторические первоисточники и отдельные работы об измерении треков и времени пролёта.

\textbf{Текст:} В 1956 году Хэнбери Браун и Твисс опубликовали лабораторную демонстрацию корреляций фототоков: Nature, том 177, страница 27, DOI 10.1038/177027a0. Работа о Сириусе опубликована в Nature, том 178, страница 1046, DOI 10.1038/1781046a0. Без полного разбора исходных данных не называем точные параметры установки или диаметр звезды.

Пионные экспериментальные результаты обсуждались уже в 1959 году. Поэтому работу Goldhaber, Goldhaber, Lee и Pais 1960 года корректнее связывать с объяснением роли бозонной симметризации в угловых корреляциях при аннигиляции антипротона с протоном, а не называть первым измерением тяжёлоионного радиуса. Размер модельного объёма уже входил в этот расчёт; связь с размером не появилась впервые только в 1970-е годы.

Работы Копылова и Подгорецкого 1970-х годов развивали связь корреляций с координатами и временами испускания. Оптическая база приёмников и относительный импульс пионов — разные экспериментальные переменные; общий принцип не делает установки одинаковыми.

Детекторные схемы опираются на работу Anderson и соавторов о TPC STAR, а совместная идентификация — на статьи Shao и соавторов и Llope и соавторов. При вопросе о времени вернуться к различию дрейфа электронов, пролёта частицы и отдельной оценки начала события.
